\section{Metodología de Medición}
Para saber si el sistema mejora el acceso hay que medir el acceso actual y
volver a medirlo durante el piloto. El equipo usa tres instrumentos, los
mismos antes y durante el piloto, para comparar con números y no con
impresiones:
\begin{itemize}
	\item \textbf{Encuesta anónima a estudiantes} (Apéndice A del documento de
	      investigación), con preguntas cerradas. Recoge lo que solo el estudiante
	      puede decir: la espera que percibe, si olvida su credencial, si tiene
	      celular y si aceptaría mostrar un QR. La muestra es de 384 estudiantes con
	      5\,\% de error o de 196 con 7\,\%.
	\item \textbf{Entrevista a vigilancia y Control Escolar} (Apéndice B), con una
	      guía de preguntas. Aporta el procedimiento real de revisión, los casos
	      excepcionales y las condiciones para tratar datos de estudiantes.
	\item \textbf{Hoja de observación del flujo en el acceso} (Apéndice C), con
	      cronómetro y conteo, durante cinco días hábiles. Mide el tiempo de revisión
	      y la fila, que son hechos observables y no se deben preguntar.
\end{itemize}

El diseño compara un antes y un después en el mismo acceso, el mismo turno,
la misma franja horaria, los mismos observadores y el mismo instrumento. No hay
grupo de comparación, porque el acceso piloto atiende a quien llega y no es
posible operar dos accesos idénticos al mismo tiempo. Este tipo de diseño se
llama cuasiexperimental \parencite{campbellstanley1963,
hernandezsampieri2014metodologia} y ya se ha usado para evaluar un piloto
basado en código QR dentro de una institución \parencite{ortiz2026qrpilot}.
Como factores externos, por ejemplo un cambio de horario escolar, podrían
alterar los tiempos sin que el prototipo tenga que ver, los resultados se
leerán como evidencia orientativa y no concluyente. También se registra el
número de puntos de revisión activos, para no atribuir al prototipo una mejora
que venga de haber abierto más filas.

El plan de acción preliminar tiene cinco pasos: el diagnóstico con los tres
instrumentos; el diseño del hardware y de la programación del lector; el
desarrollo del servidor, la base de datos y el panel; un piloto controlado sin
retirar la revisión manual; y una evaluación que repite la observación y la
compara con la línea base.
